\documentclass[12pt]{article}
\usepackage[utf8]{inputenc}
\usepackage[T1]{fontenc}
\usepackage{hyperref}
\title{Tand{\i} Language Resources for K-12 Curriculum Development}
\author{Compiled by Claude Sonnet 5.5 (claude-sonnet-5-5) for Andr\'as Kornai}
\date{}
\begin{document}
\maketitle

\section{Introduction}
Tand\textit{\i} is written with a barred i in the native spelling; it is approximated here for pdflatex. The language (ISO 639-3 \texttt{ygm}, glottocode \texttt{yago1240}), also called Tanda, is a Finisterre language of Madang Province, Papua New Guinea. Its speakers are called Yagomi. Glottolog classifies it under Trans-New Guinea (Finisterre-Huon, Warup, Degenanic). Wikipedia gives about 20 native speakers (2020). Global Recordings Network gives a population of 140, and Joshua Project gives 400 for the Yagomi people group. These figures count different things. Little material exists.

\section{Primary Educational Resources}
None found.

\section{Linguistic and Reference Materials}
\begin{itemize}
\item Wikipedia, Tand\textit{\i} language: \url{https://en.wikipedia.org/wiki/Tand%C9%A8_language}
\item Glottolog map page for \texttt{yago1240}: \url{https://glottolog.org/resource/languoid/id/yago1240.bigmap.html}
\item Crossgram language page: \url{https://crossgram.clld.org/languages/yago1240}
\item ASJP wordlist (Tanda, PNG): \url{https://asjp.clld.org/languages/TANDA_PNG}
\item ANU project ``Nen and Tonda: two languages of southern Papua New Guinea'' appeared in the results. Tonda is a different language, so do not use it: \url{https://researchportalplus.anu.edu.au/en/projects/nen-and-tonda-two-languages-of-southern-papua-new-guinea/}
\end{itemize}

\section{Religious and Community Resources}
The search summary says Joshua Project reports no Scripture translation. Joshua Project page: \url{https://joshuaproject.net/languages/ygm}. Scripture Earth also has a listing for the code, but I did not check its contents: \url{https://www.scriptureearth.org/00eng.php?iso=ygm}

\section{Digital and Audio Resources}
Global Recordings Network has a page for \texttt{ygm}. The search summary says it holds no recordings: \url{https://globalrecordings.net/en/language/ygm}

\section{Curriculum Development Context}
With about 20 speakers (a 2020 estimate), no school use and no known teaching materials, any curriculum would start with documentation of the language by and with the community. Do not confuse the language with Tandia (\texttt{tni}) or Tonda.

\end{document}